% =====================================================================
\BookChapter[C]{C: Close to the Machine}
\label{ch:c}

\begin{researchquestion}
What does a programmer gain, and risk, by controlling memory directly?
\end{researchquestion}
\OrientationNote
% =====================================================================

\begin{factsheet}{C}
First appeared & 1972; standards C89/C90, C99, C11, C17, C23\\
Designed by & Dennis Ritchie, Bell Labs (with Ken Thompson)\\
Paradigms & imperative, procedural\\
Typing & static, fairly weak (many implicit conversions)\\
Execution & compiled to native code\\
File extensions & \file{.c} (source), \file{.h} (header)\\
Compilers & GCC, Clang/LLVM, Microsoft Visual C++, many for microcontrollers\\
Used for & operating systems, embedded systems and microcontrollers, device drivers,
  firmware, real-time control, high-performance libraries, interpreters of other languages\\
Famous software & Linux and Windows kernels, the CPython interpreter, SQLite, Git, the
  firmware in most cars, appliances and medical devices, FFTW, GSL\\
\end{factsheet}

C is the language the modern computing world is built on. Your operating system, the \py
interpreter you use every day, the firmware in your washing machine and the engine control
unit of a car are all written mostly in C. C gives the programmer direct access to memory
and almost no protection from mistakes. It is small, fast and portable, and it runs on
everything from an 8-bit microcontroller costing less than a dollar to a supercomputer.

\LearningOutcomes{Compile, link and run a C program; declare variables of fixed types and predict the result of integer arithmetic; use pointers to pass data to functions; allocate and free memory; group data in a structure; explain what a buffer overflow is; and describe how C is used in embedded systems.}

\section{History and philosophy}

\begin{history}
Ken Thompson and Dennis Ritchie created the Unix operating system at Bell Labs around 1970.
At first Unix was written in assembly language, which tied it to one computer. Ritchie
designed C (a successor to Thompson's language B) so that Unix could be rewritten in a
high-level language, and by 1973 most of the Unix kernel was in C. Because the kernel was
now portable, Unix spread to many kinds of computers --- and C spread with it. The book
\emph{The C Programming Language} by Kernighan and Ritchie (1978), known as ``K\&R'',
started the tradition of the ``Hello, world'' program.
\end{history}

C's philosophy is often summarised as ``trust the programmer''. The language does exactly
what you say, quickly, and assumes you know what you are doing. It does not check array
bounds, it does not initialise variables for you, and it does not collect garbage. This
makes C extremely efficient and predictable --- important for real-time and embedded
systems --- but also makes it easy to write programs with subtle and dangerous errors.

\section{Getting started}

On Linux and macOS a C compiler (\code{gcc} or \code{clang}) is usually installed or easy
to install. On Windows, install Visual Studio (Community edition), MSYS2, or use WSL. The
basic workflow is edit $\to$ compile $\to$ run:
\codefile{C}{c/hello.c}{hello world}
\begin{lstlisting}[style=shell]
$ gcc -std=c11 -Wall -Wextra -o hello hello.c
$ ./hello
\end{lstlisting}
\outfile{c/hello.out}
Line by line:
\begin{itemize}
\item \code{\#include <stdio.h>} is a \term{preprocessor} directive. It pastes in the
  \term{header file} that declares the standard input/output functions such as
  \code{printf}.
\item \code{int main(void)} defines the function where every C program starts. It returns
  an integer to the operating system: \code{0} means success.
\item Statements end with a semicolon \code{;} and blocks are enclosed in braces
  \code{\{ \}}. Indentation is for humans only.
\item \code{printf} prints formatted text. \code{\%d} is replaced by an integer, and
  \code{\textbackslash n} is a newline.
\end{itemize}
Always compile with warnings switched on (\code{-Wall -Wextra}). The C compiler accepts
many programs that are almost certainly wrong; the warnings tell you about them.

\section{Variables and types}

Every variable must be declared with a type before use:
\begin{lstlisting}[language=C]
int    n_steps = 40;          /* integer, usually 32 bits          */
long   count = 0L;            /* at least 32 bits                  */
float  t_single = 1.5f;       /* single precision (about 7 digits) */
double pressure = 101325.0;   /* double precision (about 16 digits)*/
char   grade = 'A';           /* one character (a small integer)   */
_Bool  done = 0;              /* or 'bool' with #include <stdbool.h> */
const double g = 9.81;        /* cannot be changed                 */
\end{lstlisting}
The exact sizes of \code{int} and \code{long} depend on the platform. When the size
matters (in embedded systems, file formats and communication protocols it always does),
include \code{<stdint.h>} and use the fixed-width types \code{int8\_t}, \code{uint16\_t},
\code{int32\_t}, \code{uint64\_t} and so on.

\begin{pitfall}[Common pitfall: uninitialised variables]
In C, a local variable that is declared but not assigned contains whatever bits happened
to be in that memory --- ``garbage''. \code{double sum; for (\ldots) sum += x[i];} may give
a different wrong answer every time you run it. Always initialise: \code{double sum = 0.0;}.
\end{pitfall}

\subsection*{There is no power operator, and integer division truncates}

C has the familiar arithmetic operators \code{+ - * /} and the remainder \code{\%}, but no
power operator: use \code{pow(x, y)} from \code{<math.h>}, or simply \code{x*x} for a
square. Integer division truncates toward zero, so \code{7/2} is \code{3} and
\code{1/3*3.0} is \code{0.0}! Write \code{1.0/3.0} when you mean real division.

C introduced shorthand that many later languages copied: \code{x += 2} means
\code{x = x + 2}, and \code{i++} increases \code{i} by one. Comparisons are
\code{== != < <= > >=}; logical operators are \code{\&\&} (and), \code{||} (or) and
\code{!} (not). In C, any non-zero value counts as ``true'' and zero as ``false''.

\begin{pitfall}[Common pitfall: \texttt{=} versus \texttt{==}]
\code{if (x = 0)} \emph{assigns} 0 to \code{x} and then tests the result, which is
false, so the block never runs. You meant \code{if (x == 0)}. Compiler warnings catch
this; another reason to always use \code{-Wall}.
\end{pitfall}

The mathematical library \code{<math.h>} provides \code{sqrt}, \code{pow}, \code{exp},
\code{log}, \code{log10}, \code{sin}, \code{cos}, \code{tan}, \code{atan2}, \code{fabs},
\code{floor}, \code{ceil}, \code{fmod}, \code{hypot}, \code{tanh} and more. On Linux you
must link it explicitly with \code{-lm}.

\section{Control flow}

The control structures of C were copied by C++, Java, C\#, JavaScript and many others, so
it is worth learning them once:
\begin{lstlisting}[language=C]
if (Re < 2300.0) {
    printf("laminar\n");
} else if (Re < 4000.0) {
    printf("transitional\n");
} else {
    printf("turbulent\n");
}

/* for (initialisation; condition; update) */
for (int i = 0; i < 10; i++) {
    sum += 1.0 / ((i + 1.0) * (i + 1.0));
}

while (fabs(step) > 1e-12) {
    step = f(x) / dfdx(x);
    x -= step;
}

do {                       /* body runs at least once */
    printf("Enter a positive number: ");
    scanf("%lf", &x);
} while (x <= 0.0);

switch (n_dims) {
    case 1:  printf("line\n");    break;
    case 2:  printf("surface\n"); break;
    default: printf("volume\n");
}
\end{lstlisting}
The C \code{for} loop is more general than \py's: it is really a \code{while} loop with
an initialisation and an update written at the top. The loop \code{for (i = 0; i < n;
i++)} runs for $i = 0, 1, \ldots, n-1$, exactly like \code{for i in range(n)}.
\code{break} and \code{continue} work as in \py. In a \code{switch}, forgetting
\code{break} makes execution ``fall through'' into the next case.

\section{Functions}

A function declares the types of its parameters and of its return value
(\code{void} if it returns nothing):
\begin{lstlisting}[language=C]
double terminal_velocity(double m, double c, double g)
{
    return sqrt(m * g / c);
}
\end{lstlisting}
A function must be declared before it is used. In larger programs the declarations
(\term{prototypes}) are collected in a header file, for example \file{fluids.h},
\begin{lstlisting}[language=C]
/* fluids.h */
double reynolds(double rho, double v, double D, double mu);
double friction_factor(double Re, double rel_roughness);
\end{lstlisting}
and the definitions are in \file{fluids.c}. Other files \code{\#include "fluids.h"} and
the linker joins everything together (Figure~\ref{fig:compile}).

C passes arguments \emph{by value}: the function receives copies. To let a function
change a variable of the caller, you pass the variable's \emph{address} --- which brings us
to pointers.

\section{Pointers and memory}

A \term{pointer} is a variable that stores a memory address. Pointers are the most
powerful and the most dangerous feature of C, and they are the main reason C is both fast
and error-prone.
\begin{itemize}
\item \code{\&x} is the address of \code{x};
\item \code{double *p} declares \code{p} as a pointer to a \code{double};
\item \code{*p} is the value stored at the address in \code{p} (``dereferencing'').
\end{itemize}
\codefile{C}{c/pointers.c}{pointers and arrays}
\outfile[title={\footnotesize\sffamily\color{muted}Output (the addresses differ from run to run)}]{c/pointers.out}
The last line shows that consecutive array elements are 8 bytes apart in memory (\code{e40}
and \code{e48} in hexadecimal): a \code{double} occupies 8 bytes. In C an array is just a
block of consecutive memory, and the array name behaves as a pointer to its first element.
This is why \code{x[i]} and \code{*(x + i)} are the same thing, and why a function cannot
find out the length of an array passed to it: it receives only the starting address.

\begin{pitfall}[Common pitfall: buffer overflows]
C does not check array bounds. If \code{double a[10];} and you write \code{a[10] = 1.0;},
you overwrite whatever lies after the array in memory --- another variable, or the
function's return address. The program may crash, give wrong results, or appear to work.
Such \term{buffer overflows} are the cause of a large fraction of all security
vulnerabilities in software, which is why governments and industry now encourage
\emph{memory-safe} languages (Rust, Go, Java, C\#, \py) for new projects.
\end{pitfall}

\subsection*{Memory requested at run time must be given back}

Arrays declared as \code{double a[100];} have a size fixed when the program is compiled.
When the size is known only at run time --- for example, the number of lines in a data
file --- you request memory from the heap with \code{malloc} and release it with
\code{free}. The next program reads an unknown number of values from a file into an array
that doubles its capacity whenever it is full (exactly what a \py list does internally),
and computes statistics. It also shows a \code{struct}, C's way of grouping related values.
\codefile{C}{c/readstats.c}{reading a file into a dynamic array}
\begin{lstlisting}[style=shell]
$ gcc -std=c11 -Wall -o readstats readstats.c -lm
$ ./readstats ../data/yield.txt
\end{lstlisting}
\outfile{c/readstats.out}
Compare this with the short \py program \file{read\_csv.py} in Chapter~\ref{ch:python}, which does a similar job.
The C version is longer because it does by hand what \py does automatically: managing
memory, checking that the file exists, and growing the array. In return, it uses exactly as
much memory as needed and runs on a system with no operating system at all.

\section{Structures}

A \code{struct} groups values of different types under one name, like a \py class without
methods:
\begin{lstlisting}[language=C]
typedef struct {
    char   name[32];
    double density;          /* kg/m^3 */
    double youngs_modulus;   /* Pa     */
} Material;

Material steel = {"S275", 7850.0, 210e9};
printf("%s: E = %.0f GPa\n", steel.name, steel.youngs_modulus / 1e9);

Material *pm = &steel;
pm->density = 7860.0;        /* -> accesses a member through a pointer */
\end{lstlisting}
Structures are the C programmer's tool for building data structures: records, linked lists,
trees, and the ``objects'' that C++ later formalised.

\section{Input and output}

\code{printf} and its relatives use \term{format specifiers}: \code{\%d} (int),
\code{\%ld} (long), \code{\%zu} (size), \code{\%f}, \code{\%e} and \code{\%g} (double),
\code{\%s} (string), \code{\%c} (character), \code{\%p} (pointer). Width and precision go
between \code{\%} and the letter: \code{\%10.4f}. \py's f-string format codes come
directly from these. Reading uses \code{scanf} and \code{fscanf}; note that for a
\code{double} the reading format is \code{\%lf} but the printing format is \code{\%f}.

Files are handled through a \code{FILE *} pointer: \code{fopen}, \code{fprintf},
\code{fscanf}, \code{fgets}, \code{fclose}. Always check whether \code{fopen} returned
\code{NULL}.

\section{Case study: the falling skydiver}

\codefile{C}{c/drag.c}{Euler solution of the drag equation}
\begin{lstlisting}[style=shell]
$ gcc -std=c11 -O2 -Wall -o drag drag.c -lm
$ ./drag
\end{lstlisting}
\outfile{c/drag.out}
Again identical to \py and Fortran. The array size is a compile-time constant
(\code{\#define N\_STEPS 40}), which is typical of embedded C, where dynamic memory is often
forbidden. \code{v * v} is used instead of \code{pow(v, 2)}, a common and slightly faster
idiom.

\section{C in embedded systems}

Most C code in the world does not run on PCs but on \term{microcontrollers}: tiny
computers on a single chip that control motors, read sensors and run machines. An Arduino
board is a friendly introduction; its ``sketches'' are C++ but look almost like C. The
following program for an Arduino reads a temperature sensor and switches a fan on and off
with \term{hysteresis}, so that the fan does not chatter around the set point:
\begin{lstlisting}[language=C]
const int SENSOR_PIN = A0;
const int FAN_PIN = 9;
const float T_ON = 30.0, T_OFF = 27.0;   /* degrees C */
int fan_on = 0;

void setup() {
    pinMode(FAN_PIN, OUTPUT);
    Serial.begin(9600);
}

void loop() {                             /* called again and again, forever */
    int raw = analogRead(SENSOR_PIN);     /* 0 ... 1023 */
    float volts = raw * 5.0 / 1023.0;
    float T = (volts - 0.5) * 100.0;      /* TMP36 sensor: 10 mV per degree */
    if (!fan_on && T > T_ON)  fan_on = 1;
    if (fan_on && T < T_OFF)  fan_on = 0;
    digitalWrite(FAN_PIN, fan_on ? HIGH : LOW);
    Serial.println(T);
    delay(500);                           /* wait 0.5 s */
}
\end{lstlisting}
Embedded C has its own rules: memory is measured in kilobytes, there may be no
floating-point hardware, timing must be predictable, and the program must never crash.
Safety-critical industries use coding standards such as \textbf{MISRA C} (automotive) that
forbid the most error-prone features of the language.

\section{Strengths and weaknesses}

\begin{center}
\small
\begin{tabularx}{\linewidth}{@{}>{\raggedright\arraybackslash}X>{\raggedright\arraybackslash}X@{}}
\toprule
\textbf{Strengths} & \textbf{Weaknesses}\\
\midrule
Very fast; small and predictable executables & No protection against memory errors\\
Runs on every processor ever made & Manual memory management\\
Direct access to hardware and memory & Verbose for numerical work; no built-in arrays of variable size or matrices\\
Small language that can be learnt completely & Weak string handling\\
Stable; the common language that other languages call (the ``C ABI'') & Undefined behaviour: some errors have no defined result at all\\
\bottomrule
\end{tabularx}
\end{center}

\section{Further reading}
\begin{itemize}
\item B.~W.~Kernighan and D.~M.~Ritchie, \emph{The C Programming Language}, 2nd ed.,
  Prentice Hall, 1988.
\item K.~N.~King, \emph{C Programming: A Modern Approach}, 2nd ed., Norton, 2008.
\item J.~Gustedt, \emph{Modern C}, Manning (the current edition covers C23; a free version is online).
\item W.~H.~Press et al., \emph{Numerical Recipes in C}, 2nd ed., Cambridge University
  Press, 1992 --- the classic collection of numerical algorithms.
\end{itemize}

\begin{exercises}
\item Compile \file{drag.c} without \code{-lm}. What error do you get, and at which stage
  (compiling or linking)?
\item What do the following expressions evaluate to in C? \code{7/2}, \code{7.0/2},
  \code{7/2*2.0}, \code{(double)7/2}, \code{-7/2}, \code{-7\%2}. Check with a short program.
\item Write a C function \code{void polar(double x, double y, double *r, double *theta)}
  that returns the polar coordinates of a point through pointer arguments. Call it from
  \code{main}.
\item Write a program that reads a matrix size $n$ from the keyboard, allocates an
  $n\times n$ array of doubles with \code{malloc}, fills it with the Hilbert matrix
  $H_{ij} = 1/(i+j+1)$ (zero-based indices), prints it, and frees the memory.
\item Introduce a deliberate bug into \file{readstats.c}: remove the line
  \code{free(x);}. The program still works. Why is it nevertheless wrong? Research the
  tool \code{valgrind} (or the compiler option \code{-fsanitize=address}) and use it to
  detect the leak.
\item Write a \code{struct} for a resistor (name, resistance, power rating) and a function
  that computes the equivalent resistance of an array of resistors in series and in
  parallel.
\end{exercises}
